\documentclass[aspectratio=169]{beamer}
\usetheme{Madrid}
\usecolortheme{beaver}
\usepackage[utf8]{inputenc}
\usepackage[T1]{fontenc}
\usepackage{tikz}
\usepackage{amssymb}
\usepackage{graphicx}
\setbeamertemplate{navigation symbols}{}
\setbeamercolor{title}{fg=white,bg=cyan!60!black}
\setbeamercolor{frametitle}{fg=white,bg=cyan!50!black}
\setbeamercolor{block title}{bg=cyan!15,fg=black}
\setbeamercolor{block body}{bg=cyan!5,fg=black}

\title{G10 — Design \& Capstone}
\subtitle{Plan → build → test → ship YOUR game}
\author{ByteForge --- Grade 7 --- Game Maker --- Week Overview}
\date{}

\begin{document}

\begin{frame}
    \titlepage
\end{frame}

\begin{frame}{This Week's Plan}
    \begin{itemize}
        \item Lesson 1: Design Your Game
        \item Lesson 2: Score \& Win Screen
        \item[$\bigstar$] BOSS --- Ship Your Own Game (\texttt{lessonBoss.tex})
    \end{itemize}
\end{frame}

\begin{frame}{Real-World Example}
    \Large
    Every shipped game went through exactly this: a design document, a build, a playtest, then changes based on what real players actually did — not what the designer assumed they'd do. Minecraft itself is still playtested and balanced this way every update (mob health, drop rates, tool durability all get tuned from real player data).
\end{frame}


\begin{frame}{Let Them Discover}
    \Large
    Before opening the platform: sketch your game idea on paper — one sentence for "how do you win", and one sketch of what the screen looks like. Trade with a partner and ask: is the win condition actually checkable?
\end{frame}

\begin{frame}{Files in This Folder}
    \small
    \begin{itemize}
        \item \texttt{lesson1.tex} --- Design Your Game
        \item \texttt{lesson2.tex} --- Score \& Win Screen
        \item \texttt{lessonBoss.tex} --- BOSS: Ship Your Own Game
    \end{itemize}
    \vspace{0.5em}
    \normalsize Each lesson file is its own short deck --- open the one for
    the specific day you're teaching, not this whole overview.
\end{frame}

\end{document}
